% UTF-8; include after loading amsmath and Chinese text support.
% Generated from the Markdown report; equations and results share one source.
\section{共享只读缓存下的事件模型与有界搜索}
所有候选使用原始图与固定配置，每次评估从空缓存开始，不读取历史方案。
容量 $C_2=1048576$、$C_{\rm L1}=524288$、$C_{\rm UB}=131072$ bytes；
带宽 $B_D=60$、$B_2=250$ bytes/cycle。
$G=(V\cup\mathcal T,E)$ 为原图，$P=(\phi,\kappa,\pi)$ 为子图、分核与顺序；
$s_q,f_q$ 为原评估推导的操作开始与完成时刻，$b_x$ 为张量字节数。
$Q_e$ 为事件后的 FIFO 队列，$Q_0=\varnothing$；命中不刷新顺序。
COPY\_IN 完成时对不驻留且可容纳的张量尝试插入；COPY\_OUT 不填充缓存。

跨核读取释放约束：L2 命中不能省略 500 cycles 同步。
\[
s_{\mathrm{in}}\ge f_{\mathrm{out}}+500. \tag{1}
\]

私有存储容量：实际峰值由附录 C 的展开与换入换出确定。
\[
\max_t M_k^r(t)\le C_r,\quad r\in\{\mathrm{L1},\mathrm{UB}\}. \tag{2}
\]

发射时命中判定：查询此前状态，不刷新 FIFO 次序。
\[
h_q=\mathbf1\{x(q)\in Q_{e(q)-1}\}. \tag{3}
\]

FIFO 插入：删除能容纳新条目的最短队首前缀，已驻留时不变。
\[
j^*=\min\left\{j\in\{0,\ldots,m\}:\sum_{i=j+1}^{m}b_{x_i}+b_x\le C_2\right\},
\quad Q'=(x_{j^*+1},\ldots,x_m,x). \tag{4}
\]

驻留窗口：a 和 d 为事件编号，同周期按实际处理次序区分。
\[
\mathcal W_x=\bigcup_j[a_{xj},d_{xj}),\qquad
h_q=\mathbf1\{e(q)\in\mathcal W_{x(q)}\}. \tag{5}
\]

参考工作量：操作汇总字节按原函数取值，带宽单位 bytes/cycle。
\[
\ell_q=\max\left(1,\left\lceil b_q^{\mathrm{op}}/B_{r(q)}\right\rceil\right).
\tag{6}
\]

共享服务积分：N 为正剩余工作请求数，数值执行还须原程序的整数取整。
\[
w_q(t)=\left[\ell_q-\int_{s_q}^{t}
\frac{\mathbf1\{q\text{仍有未完成工作}\}}{N_{r(q)}(u)}\,du\right]_+ . \tag{7}
\]

主目标：完成时间最小，同时间以官方额外搬运量破同。
\[
\min_{P\in\mathcal F_n}T_2(P),\qquad T_2(P)=\max_q f_q(P). \tag{8}
\]

字节命中率：无查询字节取零，查询字节与操作汇总字节按原字段区分。
\[
H(P)=\frac{\sum_{q\in\mathrm{COPY\_IN}}b_{x(q)}h_q}
{\sum_{q\in\mathrm{COPY\_IN}}b_{x(q)}}. \tag{9}
\]

双负载分核：M 为矩阵 Pipe，O 为其他非 COPY 操作工作，这是代理量。
\[
k^*=\arg\min_k\left(
\max(L_{kM}+W_{jM},L_{kO}+W_{jO}),\
L_{kM}+L_{kO},\ k\right). \tag{10}
\]

分带割边代理：可能重复计入多消费关系，不作为精确搬运量或下界。
\[
B_{\mathrm{cut}}(\ell)=
\sum_{(u,x,v):\,d(u)\le\ell<d(v)}b_x. \tag{11}
\]

错峰排序代理：没有精确计入共享竞争与重叠，接受由完整评估决定。
\[
\delta(S)=\left|\widehat b_{\mathrm{in}}(S)/60-
(t_{\mathrm{fill}}-s_{\mathrm{miss}})\right|. \tag{12}
\]

同方案两配置对照：逐例比值算术平均，不是两边分别优化所得最优值之比。
\[
R_{i,n}=\frac{T_B(P_{i,n})}{T_2(P_{i,n})},\qquad
\overline R_{\mathcal I,n}=\frac1{|\mathcal I|}\sum_{i\in\mathcal I}R_{i,n}. \tag{13}
\]

\paragraph{算法与证据边界}
每次独立求解比较完整分量装箱、分量块顺序和数据感知深度带，
再由本次缓存事件提出至多两个合法重排候选，总计不超过五次原评估。
本轮七例、三种核数共二十一组，确认组三例五核同方案配置比值均值为 1.007435。
该值不是全量均值。四个反馈候选均未胜出，未证明局部时序修正有稳定优势。
缓存窗口经完整事件重放核对；积分表达不替代原整数事件模拟器。
代码与参数在确认前冻结，隔离改名输入复现了完整方案与结果。
